\documentclass[11pt]{report}

% ============================================================================
%  Interface-Limited Realism: A Formal Monograph
%  Full formal LaTeX. Compile with pdflatex (twice for refs).
% ============================================================================

\usepackage[a4paper,margin=1in]{geometry}
\usepackage{amsmath,amssymb,amsthm,mathtools}
\usepackage{bm}
\usepackage{enumitem}
\usepackage{microtype}
\usepackage[hidelinks]{hyperref}
\usepackage{xcolor}
\usepackage{booktabs}
\usepackage{longtable}

% --- theorem environments ---------------------------------------------------
\theoremstyle{plain}
\newtheorem{theorem}{Theorem}[chapter]
\newtheorem{proposition}[theorem]{Proposition}
\newtheorem{lemma}[theorem]{Lemma}
\newtheorem{corollary}[theorem]{Corollary}
\theoremstyle{definition}
\newtheorem{definition}[theorem]{Definition}
\newtheorem{construction}[theorem]{Construction}
\newtheorem{example}[theorem]{Example}
\theoremstyle{remark}
\newtheorem{remark}[theorem]{Remark}

% --- status tags ------------------------------------------------------------
\newcommand{\status}[1]{\par\smallskip\noindent{\footnotesize\textsc{Status:}~#1}\par}
\newcommand{\PROVEN}{\textsc{[proven]}}
\newcommand{\DERIVED}{\textsc{[derived]}}
\newcommand{\CONSTRUCTED}{\textsc{[constructed]}}
\newcommand{\STANDARD}{\textsc{[standard]}}
\newcommand{\OPENs}{\textsc{[open]}}
\newcommand{\SKETCH}{\textsc{[sketch]}}
\newcommand{\COND}{\textsc{[conditional]}}

% --- notation ---------------------------------------------------------------
\newcommand{\Real}{\mathcal{R}}
\newcommand{\Hil}{\mathcal{H}}
\newcommand{\Balg}{\mathcal{B}}
\newcommand{\Alg}{\mathcal{A}}
\newcommand{\Iint}{\mathcal{I}}
\newcommand{\Fix}{\mathrm{Fix}}
\newcommand{\Tr}{\operatorname{Tr}}
\newcommand{\dsem}{d_{\mathrm{sem}}}
\newcommand{\dgraph}{d_{\mathrm{graph}}}
\newcommand{\dembed}{d_{\mathrm{embed}}}
\newcommand{\Phiface}{\Phi_{\mathrm{iface}}}
\newcommand{\Phirec}{\Phi_{\mathrm{rec}}}
\newcommand{\PiC}{\Pi_{C}}
\newcommand{\Grep}{\Gamma_{\mathrm{rep}}}
\newcommand{\kcurv}{\kappa_{\mathrm{curv}}}
\newcommand{\kcontr}{\kappa_{\mathrm{contr}}}
\newcommand{\norm}[1]{\lVert #1 \rVert}
\newcommand{\inner}[2]{\langle #1, #2 \rangle}
\DeclareMathOperator{\Cnt}{Cnt}
\newcommand{\ket}[1]{\lvert #1 \rangle}
\newcommand{\bra}[1]{\langle #1 \rvert}
\newcommand{\braket}[1]{\langle #1 \rangle}

\title{\bfseries Interface-Limited Realism\\[2mm]
\large A Formal Framework for Observer-Constrained Knowledge\\
and the Regulation of Semantic Drift}
\author{Robert Ringler\thanks{Computational verification artifacts (machine-checked
proofs and reproducible numerics) accompany this monograph and are referenced
throughout.}}
\date{\today}

\begin{document}
\maketitle

\begin{abstract}
We develop a formal framework for observer-constrained access to reality. Reality
is represented by a Polish state space $(\Real, d_\Real)$; observation is mediated
by completely positive trace-preserving (CPTP) interface maps. We show that
observers access not underlying states but equivalence classes induced by interface
constraints, and that the structure surviving projection across a family of
heterogeneous interfaces---the \emph{interface invariant}
$\Iint=\bigcap_i \Fix(\Phi_i)$---provides a mathematically grounded notion of
objective knowledge. We prove that observation is non-expansive (the data-processing
inequality), that under a faithful invariant state the interface-invariant
observables form a von Neumann algebra carrying a conditional expectation, and that
the invariant of a family of interfaces is again such an algebra. We then construct
a metric stabilization layer: a complete semantic metric $\dsem$, an observer
closure operator $\PiC$ whose fixed-point set is the intent set, and a strictly
contractive recursive stabilizer $\Phirec$ with explicit Lipschitz constant
$\kcontr=1-\delta$, yielding Banach convergence to a unique equilibrium and a precise
drift-containment bound. We analyze the compositional (functorial) structure, and we
delimit the framework's scope by a demarcation theorem: a modeling functor that
admits a model of every domain has no falsifiable content. The framework is
epistemological rather than ontological: it does not claim that observation generates
reality, nor that reality is fundamentally informational. It concerns the limits of
access to reality and the mathematical structure of observer-dependent knowledge. We
propose empirical tests based on cross-observer reconstruction, model agreement, and
invariant preservation, and we report a machine-checked distributed substrate that
realizes agreement among observers.
\end{abstract}

\tableofcontents

% ============================================================================
\chapter*{Preliminaries: claims, scope, and notation}
\addcontentsline{toc}{chapter}{Preliminaries}

\paragraph{Central theses.} (i) Reality may contain more structure than any
individual observer can access. (ii) Observation is an interface operation. (iii)
Knowledge consists of structures that remain stable under repeated projection through
constrained interfaces.

\paragraph{Scope (epistemological, not ontological).} The framework establishes a
precise distinction between \emph{underlying state} and \emph{observable state} and
derives conditions under which knowledge is stable across heterogeneous interfaces. It
does not derive physics; it does not claim consciousness creates reality; it does not
claim reality is information. Its claim concerns the \emph{limits of access} to reality
and the mathematical structure of observer-dependent knowledge.

\paragraph{Status discipline.} Every non-trivial claim is tagged: \PROVEN{} (gap-free
proof here or in a cited verification artifact), \DERIVED{} (full proof from stated
axioms), \CONSTRUCTED{} (object exhibited and properties verified), \STANDARD{}
(correct application of a classical theorem, claimed without novelty), \COND{} (valid
under a named hypothesis), \SKETCH{} (strategy with a named gap), \OPENs{} (precise
unresolved obligation). An unlabeled assertion would be a defect.

\paragraph{Notation.} $\kcurv:\mathcal M\to\mathbb R_{\ge0}$ is a curvature field
(dimension $[\text{length}]^{-2}$); $\kcontr\in[0,1)$ is a dimensionless Lipschitz
constant; these are never identified. $\Phiface$ is an interface channel; $\Phirec$ is
the recursive stabilizer; $\PiC$ is the observer closure operator. $d_\Real$ is the
(primitive) reality metric; $\dsem$ is the (constructed) semantic metric.

% ============================================================================
\part{Foundations}

\chapter{Reality, interfaces, and the observable quotient}

\section{The reality space}
\begin{definition}[Reality space]
A \emph{reality space} is a triple $(\Real, d_\Real, \tau)$ with $(\Real,d_\Real)$ a
complete separable metric (Polish) space and $\tau$ its Borel $\sigma$-algebra.
\end{definition}
\status{\STANDARD. Every Polish space is standard Borel and supports regular
conditional probabilities (Kechris~\cite{Kechris}).}

The metric $d_\Real$ is primitive: its axioms are posited for the physical
configuration space. We emphasize that $d_\Real$ is \emph{not} a semantic metric; a
distance on meanings is constructed separately in Chapter~\ref{ch:dsem}.

\section{Interfaces}
An observer does not access $\Real$ directly but through a constrained channel. We work
in the operator-algebraic setting, which subsumes both the classical (commutative) and
quantum cases.

\begin{definition}[Interface]
Let $\Alg=\Balg(\Hil)$ be the observables of a (finite-dimensional) system with state
space $\mathcal S(\Hil)$ (density operators). An \emph{interface} is a CPTP map
$\Phi:\mathcal S(\Hil)\to\mathcal S(\Hil)$ (Schr\"odinger picture), with Heisenberg
dual $\Phi^\ast:\Alg\to\Alg$ a unital completely positive map characterized by
$\Tr(\Phi(\rho)A)=\Tr(\rho\,\Phi^\ast(A))$.
\end{definition}
\status{\STANDARD; Stinespring~\cite{Stinespring}, Kraus~\cite{Kraus}.}

\begin{definition}[Indistinguishability and the observable quotient]
For an interface $\Phi$, states $\rho,\sigma$ are \emph{$\Phi$-indistinguishable},
$\rho\sim_\Phi\sigma$, iff $\Phi(\rho)=\Phi(\sigma)$. The observer accesses the
quotient $\mathcal S(\Hil)/\!\sim_\Phi$, i.e.\ the image $\Phi(\mathcal S(\Hil))$, not
the underlying state.
\end{definition}

\section{Observation does not increase distinguishability}
The single structural fact underlying interface-limited realism is that observation can
only lose information.

\begin{theorem}[Data-processing / non-expansiveness]\label{thm:dpi}
For every CPTP map $\Phi$ and all states $\rho,\sigma$,
$\norm{\Phi(\rho)-\Phi(\sigma)}_1 \le \norm{\rho-\sigma}_1$.
\end{theorem}
\status{\PROVEN.}
\begin{proof}
Write the Hermitian traceless $\rho-\sigma=\Delta_+-\Delta_-$ in its Jordan
decomposition, with $\Delta_\pm\ge0$ of orthogonal support, so
$\norm{\rho-\sigma}_1=\Tr\Delta_++\Tr\Delta_-$. Using the variational form
$\norm{X}_1=\sup_{-I\preceq M\preceq I}\Tr(MX)$, for any admissible $M$,
\[
\Tr\!\big(M(\Phi(\Delta_+)-\Phi(\Delta_-))\big)
\le \Tr\Phi(\Delta_+)+\Tr\Phi(\Delta_-)
= \Tr\Delta_++\Tr\Delta_-=\norm{\rho-\sigma}_1,
\]
by positivity of $\Phi(\Delta_\pm)$ and trace preservation. Taking the supremum over
$M$ gives the claim.
\end{proof}

\begin{corollary}
Distinguishability of states is monotone non-increasing along any chain of interfaces.
No sequence of observations can render two states more distinguishable than they were.
\end{corollary}
\status{\DERIVED{} from Theorem~\ref{thm:dpi}. Monotonicity of the relative entropy
under CPTP maps~\cite{Lindblad75,Uhlmann77} is the entropic counterpart.}

% ============================================================================
\part{Invariants and objective knowledge}

\chapter{Interface invariants}

\section{Fixed-point structure of a single interface}
\begin{definition}[Invariant observables]
An observable $A\in\Alg$ is \emph{$\Phi$-invariant} iff $\Phi^\ast(A)=A$; equivalently
$\Tr(\Phi(\rho)A)=\Tr(\rho A)$ for all $\rho$ (the interface preserves the expectation
of $A$). Write $\Fix(\Phi^\ast)=\{A:\Phi^\ast(A)=A\}$.
\end{definition}

\begin{theorem}[Fixed-point algebra]\label{thm:fixalg}
Let $\Phi$ be CPTP with a faithful invariant state $\rho_\ast$ ($\Phi(\rho_\ast)=
\rho_\ast$, $\rho_\ast>0$). Then $\mathcal N:=\Fix(\Phi^\ast)$ is a unital
$\ast$-subalgebra of $\Alg$ (a von Neumann algebra), and there exists a conditional
expectation $E_{\mathcal N}:\Alg\to\mathcal N$ compatible with $\rho_\ast$.
\end{theorem}
\status{\STANDARD; Lindblad~\cite{Lindblad99}, Frigerio~\cite{Frigerio}, Wolf~\cite{Wolf}.}
\begin{proof}[Proof sketch]
With a faithful invariant state, $\Phi^\ast$ satisfies the Schwarz inequality with
equality exactly on $\mathcal N$ (the multiplicative domain), whence $\mathcal N$ is
closed under products and adjoints; the Accardi--Cecchini / Petz conditional expectation
onto the fixed-point algebra exists because $\rho_\ast$ is faithful~\cite{Petz86}. Full
details in the cited references.
\end{proof}

The algebra $\mathcal N$ is the \emph{decoherence-free} content of the interface: the
observables whose values the interface certifies as stable.

\section{The invariant of a family: objective knowledge}
\begin{definition}[Interface invariant]
For a family of interfaces $\{\Phi_i\}_{i\in I}$ (each with a faithful invariant state),
the \emph{interface invariant} is
\[
\Iint \;:=\; \bigcap_{i\in I}\Fix(\Phi_i^\ast).
\]
\end{definition}

\begin{theorem}[Objective structure]\label{thm:objective}
$\Iint$ is a unital $\ast$-subalgebra of $\Alg$: the largest structure simultaneously
invariant under every interface in the family. It always contains the scalars
$\mathbb C\,I$, and is non-trivial (strictly larger) iff the interfaces share a common
invariant observable beyond constants.
\end{theorem}
\status{\PROVEN{} (given Theorem~\ref{thm:fixalg}).}
\begin{proof}
Each $\Fix(\Phi_i^\ast)$ is a $\ast$-subalgebra (Theorem~\ref{thm:fixalg}); an arbitrary
intersection of $\ast$-subalgebras is a $\ast$-subalgebra (closure under product, sum,
adjoint, and the unit are each preserved under intersection). Scalars are fixed by every
unital map. Non-triviality is by definition the existence of $A\notin\mathbb C I$ with
$\Phi_i^\ast(A)=A$ for all $i$.
\end{proof}

\begin{proposition}[Conserved quantities are objective]
If $A=A^\ast$ satisfies $\Phi_i^\ast(A)=A$ for all $i$ (e.g.\ $A$ commutes with all Kraus
operators of every interface---a conserved quantity), then $A\in\Iint$.
\end{proposition}
\status{\DERIVED. Physically: quantities conserved by every admissible interface
(energy, charge) are objective in the precise sense of $\Iint$.}

\section{The commutative case: common knowledge}
\begin{proposition}[Objective knowledge as the meet of $\sigma$-algebras]
In the commutative case $\Alg=L^\infty(\Omega,\mathcal F,\mu)$, an interface is a Markov
kernel; if $\Phi_i$ is the conditional expectation onto a sub-$\sigma$-algebra
$\mathcal F_i$, then $\Fix(\Phi_i^\ast)=L^\infty(\mathcal F_i)$ and
$\Iint=L^\infty\!\big(\bigwedge_i\mathcal F_i\big)$, the functions measurable with
respect to the \emph{meet} of the observers' $\sigma$-algebras.
\end{proposition}
\status{\STANDARD. The meet $\bigwedge_i\mathcal F_i$ is exactly Aumann's
common-knowledge $\sigma$-algebra~\cite{Aumann}; agents with a common prior cannot
``agree to disagree'' on $\Iint$-measurable events.}

\begin{remark}[Physical instance]
Quantum Darwinism~\cite{Zurek,OPZ} is a concrete realization: an observable is objective
when its record proliferates redundantly into many environment fragments, i.e.\ when many
independent interfaces each fix it. ``Objective $=$ redundantly recorded across interfaces''
is the physical reading of $\Iint$.
\end{remark}

\chapter{Sharpening interfaces and the limits of self-stabilization}

A natural candidate for an interface that ``concentrates'' an observer's state onto a
sharp answer is the normalized power map. We show it does not, by itself, produce a unique
stabilized state---motivating the explicit stabilization construction of Part III.

\begin{definition}
For $\beta\in(0,1)$, the \emph{sharpening map} on $\mathcal S(\Hil)$, $\dim\Hil=d\ge2$, is
$N_\beta(\rho)=\rho^\beta/\Tr(\rho^\beta)$ (functional calculus, $0^\beta:=0$).
\end{definition}

\begin{theorem}[Fixed-point structure of the sharpening map]\label{thm:K}
For every $\beta\in(0,1)$ and $d\ge2$:
\begin{enumerate}[label=(\roman*),nosep]
\item $N_\beta$ fixes both the maximally mixed state $I/d$ and every pure state; hence
it has at least two fixed points and is \emph{not} a strict contraction in any
unitarily-invariant norm;
\item its Fr\'echet derivative at $I/d$ on the trace-zero tangent space is $\beta\cdot
\mathrm{id}$ (a local $\beta$-contraction there);
\item near any pure state it is expansive, with local ratio
$\sim\varepsilon^{\beta-1}\to\infty$.
\end{enumerate}
\end{theorem}
\status{\PROVEN. Numerically confirmed in the companion artifact (the measured local
ratio equals $\beta$ and the expansion ratio equals $\varepsilon^{\beta-1}$).}
\begin{proof}
(i) For a projector $P$ (spectrum $\{1,0,\dots\}$), $P^\beta=P$ and $\Tr P^\beta=1$, so
$N_\beta(P)=P$; for $I/d$, $(I/d)^\beta=d^{-\beta}I$ and $\Tr=d^{1-\beta}$, so
$N_\beta(I/d)=I/d$. A strict contraction on the complete space $\mathcal S(\Hil)$ has a
unique fixed point (Banach), contradicting two distinct fixed points.
(ii) By the Daleckii--Krein formula the derivative of $\rho\mapsto\rho^\beta$ at $cI$ is
the scalar $\beta c^{\beta-1}$; at $c=1/d$, normalizing and using $\Tr X=0$ yields
$DN_\beta|_{I/d}[X]=\beta X$.
(iii) For $\rho_\varepsilon=(1-\varepsilon)\ket{\psi}\bra{\psi}+\varepsilon\ket{\phi}
\bra{\phi}$ ($\braket{\psi|\phi}=0$), $\norm{\rho_\varepsilon-\ket{\psi}\bra{\psi}}_1=
2\varepsilon$ while $\norm{N_\beta(\rho_\varepsilon)-\ket{\psi}\bra{\psi}}_1=2\varepsilon^
\beta/Z$, $Z=(1-\varepsilon)^\beta+\varepsilon^\beta$; the ratio
$\varepsilon^{\beta-1}/Z\to\infty$.
\end{proof}

\begin{remark}
Theorem~\ref{thm:K} shows that ``sharpening'' alone is not a stabilization mechanism:
the maximally mixed state is attracting and pure states are repelling, so iteration drives
generic states \emph{away} from sharp answers. A genuine stabilizer must be built
differently---by damped projection (Chapter~\ref{ch:phirec}).
\end{remark}

% ============================================================================
\part{Recursive stabilization}

\chapter{A complete semantic metric}\label{ch:dsem}

When interfaces act on \emph{meanings} rather than quantum states, the relevant geometry
is a metric on a semantic space. The triangle inequality is the doubtful axiom for any
``meaning distance''; we obtain it by construction.

\begin{construction}[Fused semantic metric]
On a finite carrier $S$, let $G=(S,E,w)$ be a connected weighted graph of certified
relations with shortest-path metric $\dgraph$, and let $E:S\to\mathbb R^n$ be a Lipschitz
encoder with $\dembed(x,y)=\norm{E(x)-E(y)}_2$. For $\alpha\in(0,1]$ set
$\dsem=\alpha\,\dgraph+(1-\alpha)\,\dembed$.
\end{construction}

\begin{proposition}\label{prop:metric}
$\dsem$ is a pseudometric; a genuine metric for $\alpha>0$; and on the quotient
$S/\!\sim$ a complete metric space. Its triangle-inequality violation rate is exactly $0$.
\end{proposition}
\status{\PROVEN; the zero violation rate is verified as a unit test in the companion
benchmark.}
\begin{proof}
$\dgraph$ is a shortest-path metric and $\dembed$ is the pullback of a norm, so both
satisfy (M1)--(M3); a nonnegative combination preserves each axiom. For $\alpha>0$,
$\dsem(x,y)=0\Rightarrow\dgraph(x,y)=0\Rightarrow x=y$. Finiteness gives completeness.
\end{proof}

\begin{remark}
The construction relocates the empirical risk from \emph{structure} (triangle holds by
fiat) to \emph{faithfulness} (does $\dsem$ track human judgment of meaning?), which is an
empirical question (Chapter~\ref{ch:falsification}), not a theorem.
\status{\OPENs{} (faithfulness).}
\end{remark}

\chapter{The observer fixed point as a closure operator}

\begin{theorem}[Well-posed observer fixed point]\label{thm:pic}
Let $(L,\le)$ be a complete lattice of semantic states and $O_C,R:L\to L$ with
$\PiC:=R\circ O_C$. If $\PiC$ is monotone, extensive ($x\le\PiC x$), and idempotent,
then $\PiC$ is a closure operator; its fixed-point set $\Iint:=\Fix(\PiC)$ is a nonempty
complete lattice; and $\Grep:=\PiC$ is the unique idempotent monotone reflector onto
$\Iint$.
\end{theorem}
\status{\PROVEN{} (Tarski / closure-operator theory).}
\begin{proof}
Fixed-point sets of closure operators are closed under arbitrary meets, hence form a
complete lattice (Knaster--Tarski); the top element is closed, so $\Iint\ne\emptyset$;
idempotency makes $\Grep$ a retraction, unique among monotone idempotent retractions onto
$\Iint$.
\end{proof}

This identifies the intent set with the observer-stable set ($\Iint=\Fix(\PiC)$) and the
repair map with the closure ($\Grep=\PiC$). A concrete finite $(O_C,R)$ given by Horn-clause
forward chaining realizes the hypotheses.
\status{\CONSTRUCTED{} (toy $C$). Deriving $\PiC$ from an underlying constraint geometry is
\OPENs.}

\chapter{Recursive stabilization with explicit contraction}\label{ch:phirec}

\begin{definition}[Anchored stabilizer]
In a Hilbert realization of $S/\!\sim$, with $\Iint$ nonempty closed convex, anchor $c$,
and damping $\delta\in(0,1]$, set $\Phirec(x):=P_{\Iint}\big((1-\delta)x+\delta c\big)$,
where $P_{\Iint}$ is the metric projection onto $\Iint$.
\end{definition}

\begin{theorem}[Strict contraction and unique equilibrium]\label{thm:contraction}
$P_{\Iint}$ is firmly non-expansive, and $\Phirec$ is a strict contraction with constant
$\kcontr=1-\delta$:
$\norm{\Phirec(x)-\Phirec(y)}\le(1-\delta)\norm{x-y}$.
Hence $\Phirec$ has a unique fixed point $\sigma^\ast\in\Iint$, and
$\norm{\Phirec^{k}(x)-\sigma^\ast}\le(1-\delta)^k\norm{x-\sigma^\ast}$.
\end{theorem}
\status{\DERIVED.}
\begin{proof}
Projection onto a nonempty closed convex set in Hilbert space is firmly non-expansive: the
variational inequality $\inner{u-Pu}{w-Pu}\le0$ for all $w\in\Iint$ gives
$\norm{Pu-Pv}^2\le\inner{u-v}{Pu-Pv}\le\norm{u-v}\,\norm{Pu-Pv}$. The two arguments of
$\Phirec$ differ by $(1-\delta)(x-y)$ (the $\delta c$ cancels), so
$\norm{\Phirec(x)-\Phirec(y)}\le(1-\delta)\norm{x-y}$. Banach's theorem applies on the
complete space.
\end{proof}

\begin{remark}
$\delta$ is the anchoring strength toward the reference $c$, and is measurable as
$1-(\text{observed contraction ratio})$. The modeling commitment---that drift regulation
\emph{is} anchored projection---is empirical (\OPENs), but \emph{given} it, $\kcontr<1$ is
proved, not assumed.
\end{remark}

\chapter{Stability, drift, and equivocation}

\begin{theorem}[Convergence]
Under Theorem~\ref{thm:contraction}, $\sigma^\ast$ is the unique equilibrium and
$\dsem(\Phirec^k(x),\sigma^\ast)\le(1-\delta)^k\dsem(x,\sigma^\ast)\to0$.
\end{theorem}
\status{\DERIVED{} (Banach with constructed hypotheses).}

\begin{theorem}[Drift containment]\label{thm:drift}
With per-round perturbation $x_{k+1}=\Phirec(x_k)+e_k$, $\norm{e_k}\le\delta_{\mathrm{adv}}$,
\[
\limsup_{k\to\infty}\norm{x_k-\sigma^\ast}=\frac{\delta_{\mathrm{adv}}}{1-\kcontr},
\]
a positive constant. The system is stable to a ball of this radius; ``repaired to within
$\varepsilon$'' holds iff $\delta_{\mathrm{adv}}/(1-\kcontr)\le\varepsilon$.
\end{theorem}
\status{\PROVEN.}
\begin{proof}
$\norm{x_{k+1}-\sigma^\ast}\le\kcontr\norm{x_k-\sigma^\ast}+\delta_{\mathrm{adv}}$; summing
the geometric series gives $\le\kcontr^k\norm{x_0-\sigma^\ast}+\delta_{\mathrm{adv}}/
(1-\kcontr)$; tightness follows by aligning $e_k$ with the error direction.
\end{proof}

\begin{theorem}[Equivocation is bounded]
For two $\varepsilon$-valid outputs $y,y'$,
$\dsem(\Grep y,\Grep y')\le\dsem(y,y')\le2\varepsilon$. The equilibria are within
$2\varepsilon$, not identical.
\end{theorem}
\status{\PROVEN{} (non-expansiveness of $\Grep$ plus the triangle inequality).}

% ============================================================================
\part{Compositional structure and scope}

\chapter{Functorial structure}

\begin{theorem}[Stabilization is functorial on the Kleisli category]
The assignment $F\mapsto\Grep\circ\Phirec\circ F$ is not composition-preserving on the
base category (it would force the strict contraction $\Grep\circ\Phirec$ to be the
identity), but it is a functor on the Kleisli/Eilenberg--Moore category of the monad
$\PiC$, where the monad multiplication supplies the required structural map.
\end{theorem}
\status{\PROVEN{} (defect, by a one-line counterexample on $\mathbb R$) $+$ \CONSTRUCTED{}
(the Kleisli repair).}

\begin{theorem}[Canonical realization]
Abstraction $F_{\mathrm{sem}}$ admits a right adjoint $F_{\mathrm{ont}}$ sending a semantic
state to its \emph{canonical} (maximal-entropy) realizer; the adjunction has honest unit
and counit. Uniqueness of \emph{all} realizers is false (abstraction is many-to-one);
canonical realization is the correct statement.
\end{theorem}
\status{\PROVEN{} (defect) $+$ \CONSTRUCTED{} (canonical realizer).}

\chapter{Demarcation: the limits of the framework}

\begin{theorem}[Scope theorem]\label{thm:scope}
Regard the framework as a lax monoidal functor $U:\mathbf{ConstrDom}\to\mathbf{CPTP}$.
Then (i) $U$ outputs CPTP channels on finite Hilbert spaces---strictly less than
fundamental physics; it yields no action principle and no derivation of known physical
law. And (ii) if $U$ is total on $\mathbf{ConstrDom}$---if every domain admits a model---then
$U$ has no falsifiable content.
\end{theorem}
\status{\PROVEN.}
\begin{proof}
(i) is by inspection of the codomain. (ii): a theory is falsifiable only if some possible
observation is inconsistent with it; if every domain has a model, every observation is
accommodated, so no falsifier exists.
\end{proof}

\begin{corollary}
The framework is meaningful only when restricted to domains with a genuine data-generating
process. Breadth without restriction is a defect, not a virtue: it is precisely what would
render the framework vacuous.
\end{corollary}
\status{\PROVEN.}

% ============================================================================
\part{Verification and falsification}

\chapter{A machine-checked substrate for observer agreement}

For knowledge to be objective across observers, the observers must agree; agreement among
fault-prone distributed observers is a Byzantine consensus problem.

\begin{theorem}[Cross-observer agreement, machine-checked]
Under a lock/unlock rule, no two correct observers decide conflicting values, for $n=3f+1$.
This is discharged by exhaustive model checking (TLC), symbolic induction over unbounded
rounds (Apalache), and a parametric machine-checked proof of the quorum-intersection core
(TLAPS; $143$ obligations, $0$ failed).
\end{theorem}
\status{\PROVEN{} (machine-checked; the companion artifacts contain the proofs). A
falsifier is built in: removing the lock makes an adversarial search exhibit a violation.
Fourteen protocol-level leaves of the parametric step remain \OPENs.}

\chapter{Empirical falsification}\label{ch:falsification}

The framework makes testable commitments. It is \emph{falsified} if:
\begin{enumerate}[label=(F\arabic*),nosep]
\item interface-independent invariants cannot be identified ($\Iint$ trivial);
\item cross-interface agreement performs no better than chance;
\item stable knowledge cannot be reconstructed from preserved invariants.
\end{enumerate}

\begin{definition}[Pre-registered semantic test]
On one bounded domain, with the constructed $\dsem$ and a frozen harness, faithfulness
requires, on a held-out split: triangle-violation $=0$ (structural); rank correlation with
expert judgments $\rho\ge0.70$; drift-detection $\mathrm{AUC}\ge0.85$; conformal coverage
$\ge0.90$. Falsified if every $\alpha$ reaching $\rho\ge0.70$ destroys drift detection.
\end{definition}
\status{\CONSTRUCTED{} (harness; synthetic self-test passes and the null collapses to
chance, confirming the falsifier fires). The labeled corpus is the remaining gate, \OPENs.}

These conditions permit empirical evaluation in physical (redundant environmental records),
biological (multi-sensory invariants), and social (multi-annotator semantic invariants)
systems.

\chapter{Conclusion}

We have given a formal account of observer-constrained knowledge in which (i) observation
is an interface operation that cannot increase distinguishability; (ii) objective knowledge
is the interface invariant $\Iint=\bigcap_i\Fix(\Phi_i^\ast)$, a von Neumann algebra of
observables preserved across heterogeneous interfaces; (iii) drift toward intent is regulated
by a strictly contractive anchored projection with explicit constant $\kcontr=1-\delta$; and
(iv) the framework's scope is delimited by a demarcation theorem and grounded in a
machine-checked agreement substrate and a pre-registered falsification protocol. The
framework is epistemological: it concerns the limits of access to reality and the structure
of observer-dependent knowledge, not the generation of reality. Its central empirical
obligation---whether faithful interface invariants exist and agree across observers in real
systems---is precisely stated and testable.

% ============================================================================
\begin{thebibliography}{99}
\bibitem{Aumann} R.~J.~Aumann, \emph{Agreeing to Disagree}, Ann.\ Statist.\ \textbf{4}
(1976) 1236--1239.
\bibitem{Blackwell} D.~Blackwell, \emph{Equivalent comparisons of experiments}, Ann.\
Math.\ Statist.\ \textbf{24} (1953) 265--272.
\bibitem{Frigerio} A.~Frigerio, \emph{Stationary states of quantum dynamical semigroups},
Comm.\ Math.\ Phys.\ \textbf{63} (1978) 269--276.
\bibitem{Kechris} A.~S.~Kechris, \emph{Classical Descriptive Set Theory}, Springer, 1995.
\bibitem{Kraus} K.~Kraus, \emph{States, Effects, and Operations}, Lecture Notes in Physics
\textbf{190}, Springer, 1983.
\bibitem{Lindblad75} G.~Lindblad, \emph{Completely positive maps and entropy inequalities},
Comm.\ Math.\ Phys.\ \textbf{40} (1975) 147--151.
\bibitem{Lindblad99} G.~Lindblad, \emph{A general no-cloning theorem}, Lett.\ Math.\ Phys.\
\textbf{47} (1999) 189--196; and standard results on fixed-point sets of channels.
\bibitem{OPZ} H.~Ollivier, D.~Poulin, W.~H.~Zurek, \emph{Objective properties from subjective
quantum states: Environment as a witness}, Phys.\ Rev.\ Lett.\ \textbf{93} (2004) 220401.
\bibitem{Petz86} D.~Petz, \emph{Sufficient subalgebras and the relative entropy of states of
a von Neumann algebra}, Comm.\ Math.\ Phys.\ \textbf{105} (1986) 123--131.
\bibitem{Stinespring} W.~F.~Stinespring, \emph{Positive functions on $C^\ast$-algebras},
Proc.\ Amer.\ Math.\ Soc.\ \textbf{6} (1955) 211--216.
\bibitem{Uhlmann77} A.~Uhlmann, \emph{Relative entropy and the Wigner--Yanase--Dyson--Lieb
concavity in an interpolation theory}, Comm.\ Math.\ Phys.\ \textbf{54} (1977) 21--32.
\bibitem{Wolf} M.~M.~Wolf, \emph{Quantum Channels and Operations: Guided Tour}, lecture
notes, 2012.
\bibitem{Worrall} J.~Worrall, \emph{Structural Realism: The Best of Both Worlds?},
Dialectica \textbf{43} (1989) 99--124.
\bibitem{Zurek} W.~H.~Zurek, \emph{Quantum Darwinism}, Nature Physics \textbf{5} (2009)
181--188.
\end{thebibliography}

\end{document}
